\documentclass[]{../common_P/InfoII_P} %Put "solution" into option list to print out solutions


\usepackage{epsfig}
\usepackage{subfigure}
\usepackage{amsfonts}
\usepackage{hyperref}
\usepackage[english, ngerman]{babel}
\usepackage[utf8]{inputenc}
\usepackage[T1]{fontenc}
\usepackage{fancyvrb}
\usepackage{import}
\usepackage{eurosym}
\usepackage{wrapfig}
\def\eule{\includegraphics[width=45mm]{../common_P/unilogo.pdf}}
\SetCGPageDimensions
\thispagestyle{empty}
\usepackage{hyperref}
\hypersetup{pdftitle={},
  pdfsubject={},
  pdfauthor={},
  pdfkeywords={},
  pdfpagemode={},
  pdfborder={0 0 0 0 0 0},
  colorlinks={true}
  % linkcolor={darkblue},
  % citecolor={darkblue},
  % filecolor={black},
  % urlcolor={black}
}

\usepackage{listings}

%\lstset{numbers=left, numberstyle=\tiny, basicstyle=\tiny, language=C++}
\lstset{
  language=Scheme,
  float,
  mathescape=true,
  texcl=true,
  basicstyle=\scriptsize\ttfamily,
  columns=fixed,
  numbers=none,%left,
  stepnumber=1,
  captionpos=b,
  showstringspaces=false,
  keywordstyle=\color[rgb]{0.0, 0.0, 0.8},
  stringstyle=\color[rgb]{0.5, 0.0, 0.5},
  tabsize=4,
  commentstyle=\color[rgb]{0.133,0.545,0.133}\textit,
  morekeywords={ },
  frame=single,
  literate={ö}{{\"o}}1
           {ä}{{\"a}}1
           {ü}{{\"u}}1
}


\Professor{Prof. Dr. Martin V. Butz}
\Betreuer{
Dr. Jan Kneissler
}
\Vorlesung{Informatik II}


\Datum{Do, 17.04.2014}
\AbgabeKW{17}
\Blatt{1}
                
\newcommand{\todo}[1]{{\bf \textcolor{red}{ TODO: #1}}}

\begin{document}
\Header

\fbox{\parbox{\linewidth}{
F\"ur die Bearbeitung der \"Ubungsaufgaben
    haben Sie 90 Minuten Zeit. Bitte geben Sie Ihre L\"osungen der
    Aufgaben rechtzeitig im Vorlesungssystem unter dem Punkt
    Pr\"asenz\"ubungen\ ab. Dazu m\"ussen Sie sich mit Ihrem
    Benutzernamen und Passwort am Vorlesungssystem anmelden. Sie
    k\"onnen mehrmals innerhalb der Bearbeitungszeit Ihre L\"osungen
    abgeben, nur die letzte Version Ihrer Abgabe wird
    gespeichert. Wenn Sie mit den \"Ubungen fertig sind, bleiben Sie
    bitte an ihrem Platz sitzen und warten, bis ein Tutor ihre
    L\"osungen bewertet.  \textbf{Bei Fragen zur Rechnerbenutzung oder
      zu den \" Ubungsaufgaben wenden Sie sich bitte jederzeit an einen der TutorInnen!}
}}

%%-------------------------------------------------------
\section*{Aufgaben (insgesamt 22 Punkte)}
%%-------------------------------------------------------

{\bf \large Notiz:} Halten Sie sich beim Programmieren in Scheme an die eingeführte Konstruktionsanleitung:\\
a) Kurzbeschreibung als Kommentar (mittels: "; ...");\\
b) Signatur (der "Vertrag" mittels: "(: ...)");\\
c) Testfälle (mittels "(check-within ...)" oder "(check-expect ...)" oder ...)\\
d) Prozedur Gerüst und Rumpf\\
Benutzen Sie Hilfsprozeduren, wenn Teilprobleme gelöst werden müssen!\\

\import{../PueUebungen/Scheme/Pue01/}{pyramid-edge}

\import{../PueUebungen/Scheme/Pue02/}{bussgeld}

\end{document}

